\chapter{Factorize p = 3 Conditional Distributions to a Joint Distribution}
\label{app:A}
%%%%%%%%%%%%%%%%%%%%%%%%%%%%%%%%%%%%%%%%%%%%%%%%%%%%%%%%%%%%%%%%%%%%%%%%%%%%%%%%
We begin with the standard conditional univariate Gaussian distribution
\begin{equation}
    P(X_1|X_2 = x_2,X_3 = x_3) = \frac{1}{\sqrt{2\pi\sigma^2}} \exp\Bigg\{\frac{(X_1 - \mu_1)^2}{2\sigma^2}\Bigg\},
\end{equation}
where the mean of $X_1$ ($\mu_1$) is a function of $X_2$ and $X_3$. If $\sigma = 1$ and $(X_1 - \mu_1)^2$ we get
\begin{equation}
    P(X_1|X_2 = x_2,X_3 = x_3) = \frac{1}{\sqrt{2\pi}} \exp\Bigg\{-\frac{X_1^2 + \mu_1^2 - 2X_1\mu_1}{2}\Bigg\},
\end{equation}
and can rearrange
\begin{equation}
    P(X_1|X_2 = x_2,X_3 = x_3) = \frac{1}{\sqrt{2\pi}} \exp\Bigg\{X_1\mu_1 - \frac{X_1^2}{2} - \frac{\mu_1^2}{2}\Bigg\}.
\end{equation}

Our focus is on $\mu_1$, so we absorb $\frac{1}{\sqrt{2\pi}}$ and $-\frac{\mu_1^2}{2}$ in the log-normalizing constant $\Psi_1(\alpha,\beta,\omega)$ and let $C_1 = \frac{X_1^2}{2}$
\begin{equation}
    P(X_1|X_2 = x_2,X_3 = x_3) = \exp\{X_1\mu_1 - C_1 - \Phi_1(\alpha,\beta,\omega)\}
\end{equation}
where $\alpha,\,\beta$, and $\omega$ are the parameter vectors defining the mean $\mu_1$, which is a linear combination of the other variables $X_2,\,X_3$:
\begin{equation}
    \mu_1 = \alpha_1 + \beta_{2,1}X_2 + \beta_{3,1}X_3 + \omega_{2,3,1}X_2 X_3
\end{equation}
where $\alpha_1$ is the intercept, $\beta_{2,1},\,\beta_{3,1}$ are the parameters for the pairwise interactions with $X_2$ and $X_3$, and $\omega_{2,3,1}$ is the parameter for the 3-way interaction with $X_2 X_3$ (or, equivalently, the pairwise interaction with $X_2$ moderated by $X_3$; or the pairwise interaction with $X_3$ moderated by $X_2$).

We can similarly define $P(X_2|X_1 = x_1,X_3 = x_3)$ and $P(X_3|X_2 = x_2,X_1 = x_1)$ and factorize the three conditionals to obtain a joint distribution to get (rearranged)
\begin{equation}
    \begin{split}
        P(X_1,X_2,X_3) = \exp\Bigg\{&\sum_{i=1}^3 \alpha_1 X_i + \sum_{i=1}^3 \summ{j=1}{j \neq i}^3 \beta_{i,j} X_i X_j + \\
        &\omega_{1,2,3} X_1 X_2 X_3 - \sum_{i=1}^3 \big(C_i + \Psi_i(\alpha,\beta,\omega)\big)\Bigg\},
    \end{split}
\end{equation}
where we combine the parameters $\{\beta_{i,j},\,\beta_{j,i}\}$ and $\{\omega_{j,i,z},\,\omega_{i,j,z},\,\omega_{z,j,i}\}$ into single parameters by taking their sum, respectively. The joint distribution can be constructed analogously for any $p$.

%%%%%%%%%%%%%%%%%%%%%%%%%%%%%%%%%%%%%%%%%%%%%%%%%%%%%%%%%%%%%%%%%%%%%%%%%%%%%%%%
\chapter{Scaling variables eliminates correlations with interactions}
\label{app:B}
%%%%%%%%%%%%%%%%%%%%%%%%%%%%%%%%%%%%%%%%%%%%%%%%%%%%%%%%%%%%%%%%%%%%%%%%%%%%%%%%
For i.i.d. variables $X,Y$ with finite variances, let $Z=XY$. We show that $\rho(X,Z)=\frac{\text{cov}\{X,Z\}}{\sigma_X \sigma_Z} = 0$. Since $\sigma_X,\sigma_Z > 0,\;\text{then}\;\rho(X,Z) = 0$ iff $\text{cov}(X,Z)=0$.
\begin{equation*}
    \begin{split}
        \text{cov}(X,Z)&=\mathbb{E}\lbrack XZ \rbrack - \mathbb{E}\lbrack X \rbrack \mathbb{E} \lbrack Z \rbrack \\
        &= \mathbb{E}\lbrack X^2 Y \rbrack - \mathbb{E} \lbrack X \rbrack \mathbb{E} \lbrack XY \rbrack \\
        &= \mathbb{E} \lbrack X^2 \rbrack \mathbb{E} \lbrack Y \rbrack - \mathbb{E} \lbrack X \rbrack \mathbb{E} \lbrack X \rbrack \mathbb{E} \lbrack Y \rbrack \\
        &=\big(\mathbb{E} \lbrack X^2 \rbrack - (\mathbb{E} \lbrack X \rbrack)^2\big) \mathbb{E} \lbrack Y \rbrack \\
        &= \sigma_X \mathbb{E} \lbrack Y \rbrack,
    \end{split}
\end{equation*}
so $\rho(X,Z) = 0$ if $\mathbb{E} \lbrack Y \rbrack = 0$, and similarly, $\rho(Y,Z) = 0$ if $\mathbb{E} \lbrack X \rbrack = 0$.

%%%%%%%%%%%%%%%%%%%%%%%%%%%%%%%%%%%%%%%%%%%%%%%%%%%%%%%%%%%%%%%%%%%%%%%%%%%%%%%%
\chapter{Notes on partial correlations}
\label{app:C}
%%%%%%%%%%%%%%%%%%%%%%%%%%%%%%%%%%%%%%%%%%%%%%%%%%%%%%%%%%%%%%%%%%%%%%%%%%%%%%%%
If the number of observations is greater than the number of variables $(t > p)$, then the sample covariance matrix can be used to compute the partial correlations \cite{lauritzen1996graphical}. Let $Y_i^p$ denote the $p$-variate vector containing $i = 1,2,\ldots,t$ measurements over $t$ time points, where $\overline{Y}$ is the average over the time points, and $S$ is the sample covariance matrix
\begin{equation}
    S = \frac{1}{(t-1)} \sum_{i=1}^t (Y_i - \overline{Y})(Y_i - \overline{Y})\tran
\end{equation}
and $\Gamma = \Sigma^{-1}$ is the precision matrix, which is the inverse of $S$. The partial correlations are computed by multiplying the off-diagonal elements of $\Gamma$ with $-1$, and then dividing by the square root of the associated diagonal elements. The partial correlation between nodes $i$ and $j$, for instance, is
\begin{equation}
    -\frac{\gamma_{ij}}{\sqrt{\gamma_{ii}\gamma_{jj}}}
\end{equation}
Inverting $S$ to obtain $\Gamma$ requires that $S$ be positive definite---that is, that the rank of the space implied by $S$ is the same as its dimension $p$, which holds if $t > p$ (or $n > p$). If $n < p$ or $t < p$, however, then we cannot use $S$ directly and need to add information about the structure of $\Sigma$, which is interpreted as the true covariance matrix. The methods to compute partial correlations when $p > n$ commonly impose information about the sparsity (low number of edges) in the network. 

\paragraph{Partial Correlation by Shrinkage Estimation}
The shrinkage estimator $\hat{\Sigma}_S$ is obtained by a linear combination of the maximum likelihood estimate $S$ of the covariance matrix and a specified target matrix $T$:\footnote{(obtained via \code{pcor.shrink} in \code{R})} $\hat{\Sigma}_S = (1 - \lambda_{\scriptscriptstyle{S}})S + \lambda_{\scriptscriptstyle{S}} T$.

\paragraph{Partial Correlations by Moore-Penrose Inverse (Ridge Regression)}
$S^+ = \lim_{\lambda\to 0} (S\tran S + \lambda_r I)^{-1} S\tran$, where $I$ is the identity matrix, and $\lambda_r \geq 0$ is the regularization parameter. This can be calculated via the function \code{ginv} in \code{R}. 

\paragraph{Partial Correlations by Graphical Lasso Inverse}
The graphical lasso estimate of the inverse covariance matrix $\Sigma^{-1}$ is defined as the maximum of the penalized log-likelihood function $\log |\Sigma^{-1}| - \tr(S \Sigma^{-1}) - \lambda_l \norm{\Sigma^{-1}}_1$

%%%%%%%%%%%%%%%%%%%%%%%%%%%%%%%%%%%%%%%%%%%%%%%%%%%%%%%%%%%%%%%%%%%%%%%%%%%%%%%%
\chapter{Deviance and log-likelihood}
\label{app:D}
%%%%%%%%%%%%%%%%%%%%%%%%%%%%%%%%%%%%%%%%%%%%%%%%%%%%%%%%%%%%%%%%%%%%%%%%%%%%%%%%
Deviance is calculated as $D = 2(\L_{sat} - \L_{mod})$. Also, the \code{dev.ratio} is defined as: $D_R = 1 - \frac{D_{mod}}{D_0}$

Unrelated to this, VAR models with $p$ variables will have $p + p^2$ parameters. The Lasso-VAR has several advantages over the Bayesian VAR as it is more computationally tractable in high dimensions, performs variable selection, and can readily compute multi-step forecasts and their associated prediction intervals.

Given some model $Y = \bX\bbeta + \beps$, imagine that we obtain $\hat{\bbeta} = \pinv{\bX\tran\bX}\bX\tran Y$, along with $\hat{Y} = \bX\hat{\bbeta}$. We then assume that $\hat{\mu} = \hat{Y}$, and $\hat{\sigma}^2 = \frac{1}{N}\sum\limits_{i=1}^N (y_i - \hat{y}_i)^2$, allowing us to find the log-likelihood of the model:
\begin{equation}
    \ln{\L(\hat{\mu},\,\hat{\sigma}^2)} = \sum_{i=1}^N \ln{f(y_i\,|\,\hat{\mu},\,\hat{\sigma}^2)} = -\frac{N}{2}\ln{2\pi} - \frac{N}{2}\ln{\hat{\sigma}^2} - \frac{1}{2\pi\hat{\sigma}^2} \sum_{i=1}^N (y_i - \hat{\mu}_i)^2
\end{equation}

For binomial models, we instead specify the log likelihood as:
\begin{equation}
    \ln\L(\bbeta) = \ell(\bbeta) = \sum_{i=1}^N y_i \bigg(\sum_{k=0}^K x_{ik} \beta_{k}\bigg) - \log(1 + \exp{\sum_{k=0}^K x_{ik}\beta_k})
\end{equation}
or in matrix notation:
\begin{equation}
    \ln\L(\bbeta) = \ell(\bbeta) = \sum_{i=1}^N y_i \bx_i\bbeta - \log\big(1 + \exp(\bx_i\bbeta)\big)
\end{equation}

%%%%%%%%%%%%%%%%%%%%%%%%%%%%%%%%%%%%%%%%%%%%%%%%%%%%%%%%%%%%%%%%%%%%%%%%%%%%%%%%
\chapter{LASSO Regularization}
\label{app:E}
%%%%%%%%%%%%%%%%%%%%%%%%%%%%%%%%%%%%%%%%%%%%%%%%%%%%%%%%%%%%%%%%%%%%%%%%%%%%%%%%
There are multiple types of LASSO regularization to consider. First, take the $\ell_1$-regularization where we attempt to find $\hat{\beta}$
\begin{equation}
    \hat{\beta} = \argmin_{\beta} \lb \frac{1}{N}\norm{y - X \beta}_2^2 + \lambda\norm{\beta}_1 \rb 
\end{equation}
The above (already written in Lagrangian form) can be re-written as 
\begin{equation}
    \hat{\beta} = \argmin_{\beta} \lb \frac{1}{N}\sum_{i=1}^N (y_i - X_i \beta)^2 + \lambda\sum_{j=1}^p |\beta_j| \rb 
\end{equation}
Similarly, we have the $\ell_2$-regularization, otherwise referred to as \textit{ridge regression}:
\begin{equation}
    \hat{\beta} = \argmin_{\beta} \lb \frac{1}{N}\norm{y - X \beta}_2^2 + \lambda\norm{\beta}_2^2 \rb 
\end{equation}
which can be rewritten as:
\begin{equation}
    \hat{\beta} = \argmin_{\beta} \lb \frac{1}{N}\sum_{i=1}^N (y_i - X_i \beta)^2 + \lambda\sum_{j=1}^p \beta_j^2 \rb 
\end{equation}
Lastly, we have the \textit{elastic net}, which considers both together
\begin{equation}
        \hat{\beta} = \argmin_{\beta} \lb \frac{1}{N}\norm{y - X \beta}_2^2 + \lambda_1\norm{\beta}_1 + \lambda_2\norm{\beta}_2^2 \rb 
\end{equation}
where we are looking to know
\begin{equation}
    \min_{\beta} \lb\frac{1}{N}\sum_{i=1}^N (y_i - X_i\beta)^2 \rb \; \text{subject to} \; \frac{(1-\alpha)}{2}\sum_{j=1}^p\beta_j^2 + \alpha\sum_{j=1}^p|\beta_j| \leq t
\end{equation}
Finally, we have the \textit{adaptive lasso} which takes the form $RSS(\beta) + \lambda\sum_{j=1}^p \hat{w}_j|\beta_j|$
\begin{equation}
    \hat{\beta} = \argmin_{\beta}\lb \sum_{i=1}^N (y_i - X_i\beta)^2 + \lambda\sum_{j=1}^p \hat{w}_j |\beta_j| \rb \; \text{where} \; \hat{w}_j = \frac{1}{|\tilde{\beta}_j|^{\gamma}}, \: \forall\gamma > 0.
\end{equation}
These methods are implemented in the \code{glmnet} package within \code{R} \cite{friedman2010bregularization}, wherein generalized linear models are estimated with $\ell_1$, $\ell_2$, and the elastic net using cyclical coordinate descent, computed along a regularization path. 

%%%%%%%%%%%%%%%%%%%%%%%%%%%%%%%%%%%%%%%%%%%%%%%%%%%%%%%%%%%%%%%%%%%%%%%%%%%%%%%%
\chapter{More LASSO Stuff}
\label{app:F}
%%%%%%%%%%%%%%%%%%%%%%%%%%%%%%%%%%%%%%%%%%%%%%%%%%%%%%%%%%%%%%%%%%%%%%%%%%%%%%%%
The general form for a standard $\ell^p$ vector norm is $\norm{\beta}_p = \left(\sum_{i=1}^N|\beta_i|^p\right)^{1/p}$
\begin{equation}
    \norm{\beta}_p = \Bigg(\sum_{i=1}^N |\beta_i|^p\Bigg)^{1/p}
\end{equation}
The original form of the LASSO \cite{tibshirani1996regression} was
\begin{equation}
    (\hat{\alpha},\hat{\beta}) = \argmin_{\alpha \in \mathbb{B},\beta \in \mathbb{R}^p} \lb \sum_{i=1}^{N} \Big(y_i - \alpha - \sum_{j=1}^p \beta_j x_{ij} \Big)^2 \rb \; \text{subject to} \; \sum_{j=1}^p |\beta_j| \leq t
\end{equation}
Next, we can think about the \textit{group lasso} \cite{yuan2006model}
\begin{equation}
    \hat{\beta} = \argmin_{\beta} \lb \norm{y - \sum_{j=1}^p X_j \beta_j}_2^2 + \lambda\sum_{j=1}^p \norm{\beta}_{K_j} \rb , \; \norm{z}_{K_j} = (z^t K_j z)^{1/2}
\end{equation}
which can also be written as
\begin{equation}
    \min_{\beta \in \mathbb{R}^p} \left(\norm{\matr{y}-\sum_{j=1}^p \matr{X}_j\beta_j}_2^2 + \lambda\sum_{j=1}^p \sqrt{p_j}\norm{\beta_j}_2\right)
\end{equation}
Another way to employ the group lasso is by using \textit{bi-level selection}, where one includes both a lasso and group lasso penalty---this is referred to as the \textit{sparse group lasso} (SGL)
\begin{equation}
    Q(\beta|X,y) = \L (\beta|X,y) + \lambda_1 \sum_j \sum_k |\beta_{jk}| + \lambda_2 \sum_j \norm{\beta_j}
\end{equation}
As well as the \textit{fused lasso} \cite{tibshirani2005sparsity}
\begin{equation}
    \hat{\beta} = \argmin_{\beta} \lb\frac{1}{N}\sum_{i=1}^N (y_i - x_i^t\beta)^2 \rb \; \text{subject to} \; \sum_{j=1}^p |\beta_j| \leq t_1 \; \text{and} \; \sum_{j=2}^p |\beta_j - \beta_{j-1}| \leq t_2
\end{equation}

\section{Estimating Model Parameters with the LASSO}
The lasso is not a linear fitting method, and thus there is no exact form closed solution to Cov($\matr{y},\hat{\matr{y}}$). For degrees of freedom, seems natural to say that df$(\lambda) = \norm{\hat{\beta}(\lambda)}_0$. This is allegedly an unbiased estimate of the lasso degrees of freedom.

%%%%%%%%%%%%%%%%%%%%%%%%%%%%%%%%%%%%%%%%%%%%%%%%%%%%%%%%%%%%%%%%%%%%%%%%%%%%%%%%
\chapter{Algorithm}
\label{app:G}
%%%%%%%%%%%%%%%%%%%%%%%%%%%%%%%%%%%%%%%%%%%%%%%%%%%%%%%%%%%%%%%%%%%%%%%%%%%%%%%%
\begin{spacing}{1}
\begin{algorithm}
\caption{Example table for algorithms}
\begin{algorithmic}[1]
\Procedure{Roy}{$a,b$}       \Comment{This is a test}
    \State System Initialization
    \State Read the value 
    \If{$condition = True$}
        \State Do this
        \If{$Condition \geq 1$}
        \State Do that
        \ElsIf{$Condition \neq 5$}
        \State Do another
        \State Do that as well
        \Else
        \State Do otherwise
        \EndIf
    \EndIf
    \While{$something \not= 0$}  \Comment{put some comments here}
        \State $var1 \leftarrow var2$  \Comment{another comment}
        \State $var3 \leftarrow var4$
    \EndWhile  \label{roy's loop}
\EndProcedure
\end{algorithmic}
\end{algorithm}
\end{spacing}

\begin{table}[ht]
\caption{This is a table.}
\centering
  \begin{tabular}{ccccccc}
    \toprule
    One & Two & Three & Four & Five & Six & Seven \\
    \midrule
    100 & 200 & 300 & 400 & 500 & 600 & 700 \\
    100 & 200 & 300 & 400 & 500 & 600 & 700 \\
    100 & 200 & 300 & 400 & 500 & 600 & 700 \\
    100 & 200 & 300 & 400 & 500 & 600 & 700 \\
    100 & 200 & 300 & 400 & 500 & 600 & 700 \\
    \bottomrule
  \end{tabular}
\end{table}

\begin{table}
\centering
\begin{tabular}{l|r}
Item & Quantity \\\hline
Widgets & 42 \\
Gadgets & 13
\end{tabular}
\caption{\label{tab:widgets}An example table.}
\end{table}

%%%%%%%%%%%%%%%%%%%%%%%%%%%%%%%%%%%%%%%%%%%%%%%%%%%%%%%%%%%%%%%%%%%%%%%%%%%%%%%%
\chapter{Cyclic Coordinate Descent}
\label{app:H}
%%%%%%%%%%%%%%%%%%%%%%%%%%%%%%%%%%%%%%%%%%%%%%%%%%%%%%%%%%%%%%%%%%%%%%%%%%%%%%%%
This is regular scoring algorithm that produces maximum likelihood estimates
\begin{equation}
    \begin{split}
        \nabla L(\theta) &= \sum_{i=1}^n (y_i - \hat{y}_i)x_i \\
        -d^2 L(\theta) &= \sum_{i=1}^n \hat{y}_i (1 - \hat{y}_i)x_i x_i^t
    \end{split}
\end{equation}

The score equations of the log-likelihood define part of the Karush-Kuhn-Tucker (KKT) conditions for optimality in the penalized model
\begin{align}
    \Big|\sum_{i=1}^n (y_i - \hat{y}_i^\lambda) x_{ij} \Big| = \lambda \quad \text{if} \quad \beta_j \neq 0 \\
    \Big|\sum_{i=1}^n (y_i - \hat{y}_i^\lambda) x_{ij} \Big| \leq \lambda \quad \text{if} \quad \beta_j = 0
\end{align}
and the initial absolute score would be
\begin{equation}
    a_j = \Big|\sum_{i=1}^n (y_i - \hat{y}_0) x_{ij} \Big|
\end{equation}

%%%%%%%%%%%%%%%%%%%%%%%%%%%%%%%%%%%%%%%%%%%%%%%%%%%%%%%%%%%%%%%%%%%%%%%%%%%%%%%%
\chapter{VAR and Friends}
\label{app:I}
%%%%%%%%%%%%%%%%%%%%%%%%%%%%%%%%%%%%%%%%%%%%%%%%%%%%%%%%%%%%%%%%%%%%%%%%%%%%%%%%
In Vector Autoregressive (VAR) models, the variables $\matr{X}_t \in \mathbb{R}^p$ at time point $t \in \mathbb{Z}$ are modeled as a linear combination of the same variables at earlier time points. For a lag 1 model, VAR(1), we model $\matr{X}_t$ as
\begin{equation}
    \matr{X}_t = \bm{\beta}_0 + \matr{B}\matr{X}_{t-1} + \bm{\epsilon}_t =
    \begin{bmatrix}
    X_{t,1} \\
    \vdots \\
    X_{t,p}
    \end{bmatrix} =
    \begin{bmatrix}
    \beta_{0,1} \\
    \vdots \\
    \beta_{0,p}
    \end{bmatrix} +
    \begin{bmatrix}
    \beta_{1,1} & \cdots & \beta_{1,p} \\
    \vdots & \ddots & \vdots \\
    \beta_{p,1} & \cdots & \beta_{p,p}
    \end{bmatrix}
    \begin{bmatrix}
    X_{t-1,1} \\
    \vdots \\
    X_{t-1,p}
    \end{bmatrix} + 
    \begin{bmatrix}
    \epsilon_{t,1} \\
    \vdots \\
    \epsilon_{t,p}
    \end{bmatrix}
\end{equation}

%%%%%%%%%%%%%%%%%%%%%%%%%%%%%%%%%%%%%%%%%%%%%%%%%%%%%%%%%%%%%%%%%%%%%%%%%%%%%%%%
\chapter{F change}
\label{app:J}
%%%%%%%%%%%%%%%%%%%%%%%%%%%%%%%%%%%%%%%%%%%%%%%%%%%%%%%%%%%%%%%%%%%%%%%%%%%%%%%%
Comparing two models with an $F$-test entails that: (1) There is a `restricted` model (with more residual degrees of freedom), and (2) There is a `full` model (with fewer residual degrees of freedom). The idea is to test the null hypothesis that the restricted model is sufficient to account for the variation in $Y$. The alternative hypothesis is that the full model better accounts for the variation in $Y$. 
\begin{equation}
    \begin{split}
        SS_{0} &= \sum_{i=1}^N (\hat{y}_{0,i} - y_i)^2, \quad\text{and}\quad Res\ df_0 = N - p_0 - 1 \\
        SS_{1} &= \sum_{i=1}^N (\hat{y}_{1,i} - y_i)^2, \quad\text{and}\quad Res\ df_1 = N - p_1 - 1
    \end{split}
\end{equation}
Given that $SS_0$ is taken to be the restricted model, and $SS_1$ the full model, it is true that $df_0 > df_1$ because $p_1 > p_0$.

The $F_{\Delta}$ statistic can then be calculated as follows:
\begin{equation}
    F_{\Delta} = \frac{(SS_1 - SS_0)/(df_1 - df_0)}{SS_1/df_1}, \quad\text{or}\quad F_{\Delta} = \frac{SS_{\Delta}/df_{\Delta}}{SS_1/df_1}
\end{equation}
We then find the $p$-value by indexing $F_{\Delta}$ with $df_{\Delta}$ and $df_1$: $F_{\Delta}(df_{\Delta},\,df_1)$.

For a single model, we can think of the $F$ statistic slightly differently:
\begin{equation}
    \begin{split}
        SS_0 &= \sum_{i=1}^N (y_i - \bar{y})^2 \\
        SS_1 &= \sum_{i=1}^N (y_i - \hat{y}_i)^2
    \end{split}
\end{equation}
where the residual $df$ is: $df_0 = N - p - 1$, and the standard $df$ is simply the number of predictors: $df_1 = p$. The $F$ statistic is then:
\begin{equation}
    F = \frac{(1 - \frac{SS_1}{SS_0})/df_1}{\frac{SS_1}{SS_0}/df_0}
\end{equation}
where we find the $p$-value using: $F(df_1,\,df_0)$.

%%%%%%%%%%%%%%%%%%%%%%%%%%%%%%%%%%%%%%%%%%%%%%%%%%%%%%%%%%%%%%%%%%%%%%%%%%%%%%%%
\chapter{Basics}
\label{app:K}
%%%%%%%%%%%%%%%%%%%%%%%%%%%%%%%%%%%%%%%%%%%%%%%%%%%%%%%%%%%%%%%%%%%%%%%%%%%%%%%%
\begin{gather}
    Y = \bX\bbeta + \beps, \quad\quad \hat{\bbeta} = \pinv{\bX\tran\bX}\bX\tran Y, \quad\quad \hat{Y} = \bX\hat{\bbeta} \\
    \hat{\sigma}^2 = \frac{1}{N-p-1} \sum_{i=1}^N (y_i - \hat{y}_i)^2, \quad\quad \hat{\sigma}_{\hat{\bbeta}} = \sqrt{\hat{\sigma}^2 \pinv{\bX\tran\bX}} \\
    SS_B = \sum_{i=1}^N (\hat{y}_i - \bar{\hat{y}})^2, \quad\quad SS_W = \sum_{i=1}^N (y_i - \hat{y}_i)^2, \quad\quad SS_{Total} = \sum_{i=1}^2 (y_i - \bar{y})^2 \\
    F = \frac{SS_B/p}{SS_W/(N-p-1)}
\end{gather}
If you are conducting an ANOVA with a predictor that has $k$ groups:
\begin{gather}
    RMSSE = d = \sqrt{\frac{\sum_{j=1}^k d_j^2}{k - 1}}, \quad\text{where}\quad d_j = \frac{\mu_j - \mu^*}{\sqrt{\sum_{i=1}^N (y_i - \hat{y}_i)^2 /(N - k)}} = \frac{\mu_j - \mu^*}{\sqrt{\hat{\sigma}^2}}
\end{gather}
Importantly, each $\mu_j$ is the mean of $Y$ for group $j$, and $\mu^*$ is the \textit{mean of the group means}: $\mu^* = \frac{1}{k}\sum_{j=1}^k \mu_j$. Hodges $g$ is simply:
\begin{equation}
    g = \frac{|\bar{x_1} - \bar{x_2}|}{\sqrt{MS_W}}
\end{equation}
given that $\sqrt{MS_W}$ approximates the pooled standard deviation across groups. Also, the standard error of the mean estimate for each group can be defined as: $\sqrt{MS_W/n_j}$.
\begin{equation}
    \omega^2 = \frac{SS_B - df_B MS_W}{SS_{Total} + MS_W} = \frac{SS_{Total} - df_{Total} MS_W}{SS_{Total} + MS_W}, \quad \omega_p^2 = \frac{SS_p - df_p MS_W}{SS_p + (N - df_p) MS_W}
\end{equation}
The first statement of equivalence is only true with a single predictor. Otherwise, $\omega^2$ (and $\omega_p^2$) could be calculated for each predictor individually. Another effect size, which compares one group to the reference group: $r = \sqrt{\frac{t^2}{t^2 + df_{Res}}}$. Here's how to back-solve for the SS for each predictor in a regression model: $SS_p = t_p^2 \hat{\sigma}^2$. Then, you can get effect sizes such as $\eta_{partial}^2$, where $\eta_p^2(p) = SS_p/(SS_p + SS_W)$.

When it comes to effect sizes, though, \cite{cohen1988power} agrees that standardized guidelines are not really the way to go. See p.25, and the wikipedia entry for effect size.

%%%%%%%%%%%%%%%%%%%%%%%%%%%%%%%%%%%%%%%%%%%%%%%%%%%%%%%%%%%%%%%%%%%%%%%%%%%%%%%%
\chapter{Sum of Squares}
\label{app:L}
%%%%%%%%%%%%%%%%%%%%%%%%%%%%%%%%%%%%%%%%%%%%%%%%%%%%%%%%%%%%%%%%%%%%%%%%%%%%%%%%
\begin{enumerate}
    \item Type I: Completely sequential---gives SS for a model with only the first term, then added SS for a model that adds the second term, etc.
    \item Type II: Gives independent SS for each term; however, terms that are included in an interaction are treated as though they come from a model without the interaction. That is, if $X$ is involved in an interaction, then the SS for $X$ will represent the independent SS accounted for by $X$ in a model without the interaction. And if $X$ is not involved in an interaction, the SS will be the independent SS for the model under consideration. The SS for interaction terms is treated as independent.
    \item Type III: All terms are treated as independent. The SS for any given term is proportional to the $F_{\Delta}$ of adding that term to a reduced model.
\end{enumerate}

Also, helmert coding is what SPSS uses in linear models with categorical variables. This means that the dummy coded variables are specified differently. In a `treatment` contrast, each level is compared with the reference group. In helmert coding, Level 2 is compared with Level 1; Level 3 is compared with the mean of Levels 1 and 2; Level 4 is compared with the mean of Levels 1, 2, and 3, etc. The first dummy code is created such that Level 1 is coded with $-1$, Level 2 is coded with $1$, and all other levels are coded with $0$. The second dummy code has Level 1 and 2 coded as $-1$, Level 3 coded as $2$, and all other levels coded as $0$. The third dummy code has Levels 1, 2, and 3 coded as $-1$, Level 4 coded as $3$, and all other levels coded as $0$, etc.

%%%%%%%%%%%%%%%%%%%%%%%%%%%%%%%%%%%%%%%%%%%%%%%%%%%%%%%%%%%%%%%%%%%%%%%%%%%%%%%%
\chapter{Stability Selection}
\label{app:M}
%%%%%%%%%%%%%%%%%%%%%%%%%%%%%%%%%%%%%%%%%%%%%%%%%%%%%%%%%%%%%%%%%%%%%%%%%%%%%%%%
The goal of variable selection is to identify a set of variables that will have nonzero weight in the model. We estimate the model parameters $\hat{\beta}$ and then define the selection set $\hat{S}$ as:
\begin{equation}
    \begin{split}
        \hat{\beta}^{\lambda} &= \argmin_{\beta\in\R^p} \norm{y - X\beta}_2^2 + \lambda\norm{\beta}_1 \\
        \hat{S}^{\lambda} &= \set{k:\hat{\beta}_k^{\lambda} \neq 0}
    \end{split}
\end{equation}
where $X \in \R^{n\times p}$ is the design matrix, $y \in \R^n$ is the vector of outputs, and $\lambda$ is the regularization parameter that controls the size of the selection set.

Network estimation is often formalized as an inverse covariance estimation problem in a GGM. Assuming data drawn from a multivariate Gaussian distribution $(X_1,\ldots,X_p)\sim\N(0,\Sigma)$, we can estimate the pairwise conditional dependencies among the variables $1,\ldots,p$ by estimating the inverse covariance matrix $\Theta = \Sigma\inv$ and then identify the entries with nonzero value. In order to identify some conditional independencies, we must obtain a sparse estimate of the covariance matrix. In this case, we model parameters $\hat{\Theta}$ and then define the selection set $\hat{S}$ as follows:
\begin{equation}
    \begin{split}
        \hat{\Theta}^{\lambda} &= \argmin_{\Theta\in\R^{p\times p},\Theta\succeq 0} -\log|\Theta| + \tr(S\Theta) + \lambda\norm{\Theta}_1 \\
        \hat{S}^{\lambda} &= \set{(j,k):\hat{\theta}_{j,k}^{\lambda} \neq 0}
    \end{split}
\end{equation}
where $S = \frac{1}{n}X\tran X \in \R^{p\times p}$ is the empirical covariance matrix and $\lambda$ controls the size of the selection set.

Assume we have a generic structure estimation algorithm that takes a dataset $\matr{Z} = Z_1,\ldots,Z_n$ and a regularization parameter $\lambda$ that returns a selection set $\hat{S}^{\lambda}$. We can think of this algorithm as a black box. The stability selection algorithm runs as follows:
\begin{enumerate}
    \item Define a candidate set of parameters $\Lambda$ and a subsample number $N$.
    \item For each value of $\lambda\in\Lambda$, do:
    \begin{enumerate}
        \item Start with the full dataset $\matr{Z}_{(\text{full})} = Z_1,\ldots,Z_n$
        \item For each $i$ in $1,\ldots,N$, do:
        \begin{enumerate}
            \item Subsample from $\matr{Z}_{(\text{full})}$ without replacement to generate a smaller dataset of size $\lfloor n/2 \rfloor$, given by $\matr{Z}_{(i)}$.
            \item Run the selection algorithm on dataset $\matr{Z}_{(i)}$ with parameter $\lambda$ to obtain a selection set $\hat{S}_{(i)}^{\lambda}$.
            \item Given the selection sets from each subsample, calculate the empirical selection probability for each component.
            \begin{equation}
                \hat{\Pi}_k^{\lambda} = \mathbb{P}\set{k\in\hat{S}^{\lambda}} = \frac{1}{N}\sum_{i=1}^N\mathbb{I}\set{k\in\hat{S}_i^{\lambda}}
            \end{equation}
        \end{enumerate}
    \end{enumerate}
    \item Given the selection probabilities for each component and for each value of $\lambda$, construct the stable set as:
    \begin{equation}
        \hat{S}^{\text{stable}} = \set{k:\max_{\lambda\in\Lambda}\hat{\Pi}_k^{\lambda} \geq \pi_{\text{thr}}}
    \end{equation}
    where $\pi_{\text{thr}}$ is a predefined threshold.
\end{enumerate}

An alternative to this subsampling approach is the sample splitting approach. We begin by randomly splitting the data into two non-overlapping subsamples of equal size $\lfloor n/2 \rfloor$ and see whether variable $p$ is chosen in both sets simultaneously. Let $I_1$ and $I_2$ two random subsets of $\set{1,\ldots,n}$, with $|I_i| = \lfloor n/2 \rfloor$ for $i = 1,2$ and $I_1 \cap I_2 = 0$. Define the simultaneously selected set as the intersection of $\hat{S}^{\lambda}(I_1)$ and $\hat{S}^{\lambda}(I_2)$:
\begin{equation}
    \hat{S}^{\text{simult},\lambda} = \hat{S}^{\lambda}(I_1) \cap \hat{S}^{\lambda}(I_2)
\end{equation}
with the simultaneous selection probabilities $\hat{\Pi}$ for any set $K \subseteq \set{1,\ldots,p}$ defined as:
\begin{equation}
    \begin{split}
        \hat{\Pi}_K^{\text{simult},\lambda} &= P^*(K\subseteq\hat{S}^{\text{simult},\lambda}) \\
        \hat{\Pi}_K^{\text{simult},\lambda} &\geq 2\hat{\Pi}_K^{\lambda} - 1
    \end{split}
\end{equation}
\begin{equation}
\end{equation}

%%%%%%%%%%%%%%%%%%%%%%%%%%%%%%%%%%%%%%%%%%%%%%%%%%%%%%%%%%%%%%%%%%%%%%%%%%%%%%%%
\chapter{GGM}
\label{app:N}
%%%%%%%%%%%%%%%%%%%%%%%%%%%%%%%%%%%%%%%%%%%%%%%%%%%%%%%%%%%%%%%%%%%%%%%%%%%%%%%%
Finally I understand the GGM notation.
\begin{equation}
    \bSigma = \bDelta (\I - \bOmega)\inv \bDelta
\end{equation}
where $\bSigma$ is the variance-covariance matrix, and $\bOmega$ is a matrix of partial correlations with zeros in the diagonal. Then, $\bDelta$ is a diagonal scaling matrix that "controls the variances", meaning that a diagonal element of $\bDelta$ would be $\text{diag}(\bDelta) = \text{diag}(\bSigma\inv)^{-\frac{1}{2}}$ 

Conversely, we can define the partial correlation matrix $\bOmega$ as:
\begin{equation}
    \bOmega = \I - \big(\text{diag}(\bSigma\inv)^{-\frac{1}{2}} \bSigma\inv\big)\tran \text{diag}(\bSigma\inv)^{-\frac{1}{2}}
\end{equation}
or instead:
\begin{equation}
    \bOmega = \I - \bDelta \bSigma\inv \bDelta
\end{equation}
where $\bDelta = \text{diag}(\bSigma\inv)^{-\frac{1}{2}}$.